%% LyX 2.4.4 created this file.  For more info, see https://www.lyx.org/.
%% Do not edit unless you really know what you are doing.
\documentclass[english]{article}
\usepackage[T1]{fontenc}
\usepackage[latin9]{inputenc}
\usepackage{enumitem}
\usepackage{amsmath}
\usepackage{amsthm}
\usepackage{cancel}
\usepackage{geometry}
\geometry{verbose,tmargin=1in,bmargin=1in,lmargin=1in,rmargin=1in}
\PassOptionsToPackage{normalem}{ulem}
\usepackage{ulem}

\makeatletter
%%%%%%%%%%%%%%%%%%%%%%%%%%%%%% Textclass specific LaTeX commands.
\numberwithin{equation}{section}
\numberwithin{figure}{section}
\newlength{\lyxlabelwidth}      % auxiliary length 

%%%%%%%%%%%%%%%%%%%%%%%%%%%%%% User specified LaTeX commands.
\usepackage{fancyhdr}
\usepackage{lastpage}
\pagestyle{fancy}
\usepackage{enumitem}
\usepackage{ifthen}
%sets math fonts to be more readable
\everymath=\expandafter{\the\everymath\displaystyle}
%\DeclareMathSizes{10}{10}{10}{10}
\usepackage{multicol}

\usepackage{tikz}
\usetikzlibrary{plotmarks}
\usepackage{pgfplots}
\pgfplotsset{compat=newest}

\usepackage{calc}
\usepackage[nomessages]{fp}

\usepackage{advdate}    % Advancing/saving dates
\usepackage{datetime}   % Dates formatting
\usepackage{datenumber} % Counters for dates
% Added by lyx2lyx
\setlength{\parskip}{\smallskipamount}
\setlength{\parindent}{0pt}

\makeatother

\usepackage{babel}
\begin{document}
\renewcommand{\author}{Dr.\,James Ashe}

\newcommand{\classNum}{137}

\newcommand{\classSec}{B}

\newcommand{\nameType}{Name:}

\newcommand{\examType}{Example Excercises}

\renewcommand{\date}{\formatdate{22}{3}{2013}}

%%First page's header%%%%%%%%%%%%%%%%%%%%%%
\fancypagestyle{firststyle} %%header settings for first pagr
{
	\fancyhead[L]{MTH \classNum{}(\classSec{}) \author{}\\
	\textbf{\examType{}} \date{}}
	\fancyhead[C]{\textbf{\nameType}}
	\setlength{\headsep}{14pt}
}
%%%%%%%%%%%%%%%%%%%%%%%%%%%%%%%%%%%%%%%%%%

%%Next page's header%%%%%%%%%%%%%%%%%%%%%%%
\setlist{nolistsep}
\lhead{\classNum{}(\classSec{})} 
\chead{\textbf{\nameType}} 
\cfoot{\thepage{}~of~\pageref{LastPage}} 
\setlength{\headsep}{5pt} 
%%%%%%%%%%%%%%%%%%%%%%%%%%%%%%%%%%%%%%%%%%%

%%Various settings; see the preamble for other settings%%%%%
%%\setenumerate[0]{leftmargin=12pt,itemindent=0pt,labelwidth=0pt}
\renewcommand{\labelenumi}{\textbf{\arabic{enumi}.}}
\renewcommand{\labelenumii}{\textbf{\alph{enumii}.}}
\thispagestyle{firststyle}
%%%%%%%%%%%%%%%%%%%%%%%%%%%%%%%%%%%%%%%%%%
\newcommand{\xmax}{10}
\newcommand{\xmin}{-10}
\newcommand{\xstep}{0} %must be xmin+n
\newcommand{\ymax}{10}
\newcommand{\ymin}{-10}
\newcommand{\ystep}{0} %must be ymin+n

\newcommand{\Dspace}{\vspace{1cm}}

\small

\subsubsection*{An example of solving with radicals. }
\begin{enumerate}[resume]
\item Find all real solutions of $\sqrt{x+10}+\sqrt{x-9}=19$. 

\begin{enumerate}[resume]
\item []This problem is probably easiest to tackle by isolating one of
the radicals, squaring and simplifying; and then isolating the other
radical.
\begin{eqnarray*}
\sqrt{x+10}+\sqrt{x-9} & = & 19\\
\Rightarrow\sqrt{x+10} & = & 19-\sqrt{x-9}\\
\Rightarrow\left(\sqrt{x+10}\right)^{2} & = & \left(19-\sqrt{x-9}\right)^{2}\\
\Rightarrow x+10 & = & 361-38\sqrt{x-9}+x-9\\
\Rightarrow-342 & = & -38\sqrt{x-9}\\
\left(-342\right)^{2} & = & \left(-38\sqrt{x-9}\right)^{2}\\
116964 & = & 1444\left(x-9\right)\\
81 & = & x-9\\
x & = & 90
\end{eqnarray*}
\\
While a little too complicated for an exam, this problem is good algebra
practice. The concept of isolating the radical so it can be squared
is key. 
\end{enumerate}
\end{enumerate}

\subsubsection*{Problem 12 from section 3.4: The William, Nancy, and Edgar distance
problem.}

\textbf{If they all start at the same time:}
\begin{enumerate}[resume]
\item William, Nancy, and Edgar met at the lodge. \uline{At 8 AM} William
began hiking at 2 mph, Nancy began hiking north at 5 mph, and Edgar
went East on a three-wheeler at 12 mph. At what time was the distance
between Nancy and Edgar 20 miles greater than than the distance between
Nancy and William?\vspace{.5cm}

We need to solve the inequality 
\begin{equation}
\mbox{dist}\left(N,E\right)>\mbox{dist}\left(N,W\right)+20\label{eq:dist-1}
\end{equation}
Where $N=\mbox{Nancy's position}$, $E=\mbox{Edgar's position}$,
and $W=\mbox{William's position}$ as $\left(x,y\right)$ coordinates
with the cabin at the origin (point $(0,0)$). See the figure:

\begin{center}
\renewcommand{\xmax}{15}
\renewcommand{\xmin}{-5}
\renewcommand{\xstep}{0} %must be xmin+n
\renewcommand{\ymax}{8}
\renewcommand{\ymin}{0}
\renewcommand{\ystep}{0} %must be ymin+n
%%%%%%%
\begin{tikzpicture} 
\begin{axis}[height=3.5cm, 
	width=3.5cm, 
	axis lines=middle,
	minor x tick num={4},
	minor y tick num={4},
	grid=none,
	xtick={0,50,...,250},
	ytick={0,10,20,...,40},
	%axis line style={-latex}, 
	xmin=\xmin,xmax=\xmax, 
	ymin=\ymin,ymax=\ymax,
	%restrict y to domain=-1:10,
	%no markers, 
	xlabel={$$},ylabel={$$}, 
	%every axis x label/.append style={anchor=north}, 
	%every axis y label/.append style={anchor=east}
	]
	
	\addplot[color=red,mark=o,solid] coordinates {(0,5)} node[right] {$N$};
	\addplot[color=blue,mark=o,solid] coordinates {(12,0)} node[above] {$E$};
	\addplot[color=black,mark=o,solid] coordinates {(-2,0)} node[above] {$W$};

	\addplot[color=blue,dashed,mark=*,mark size=1pt] coordinates {(0,5) (12,0) (-2,0) };
	\addplot[color=red,dashed,mark=*,mark size=1pt] coordinates {(0,5) (-2,0) };

	
\end{axis}
\end{tikzpicture} 
\par\end{center}

$N=\left(0,5t\right)$, $E=\left(12t,0\right)$, and $W=\left(-2t,0\right)$.
The $"-"$ for William nearly denotes his direction (it is not really
necessary). Each point is given with respect to time \uline{after
8 AM}. At $t=0$, they are all at the cabin; $\left(0,0\right)$.
They all move away from the cabin according to their speed and direction.
The picture shows $t=1$.

$\mbox{dist}\left(\left(x_{1},y_{1}\right),\left(x_{2},y_{2}\right)\right)=\sqrt{\left(x_{1}-x_{2}\right)^{2}+\left(y_{1}-y_{2}\right)^{2}}$,
so:
\begin{itemize}
\item $\mbox{dist}\left(N,W\right)=\sqrt{\left(0-\left(-2t\right)\right)^{2}+\left(5t-0\right)^{2}}=\sqrt{4t^{2}+25t^{2}}=\sqrt{29t^{2}}=\sqrt{29}t$
\item $\mbox{dist}\left(N,E\right)=\sqrt{\left(0-12t\right)^{2}+\left(5t-0\right)^{2}}=\sqrt{144t^{2}+25t^{2}}=\sqrt{169t^{2}}=13t$
\end{itemize}
Plugging this into equation \ref{eq:dist-1}, we have:
\begin{eqnarray*}
13t & > & \sqrt{29}t+20\\
13t-\sqrt{29}t & > & 20\\
\left(13-\sqrt{29}\right)t & > & 20\\
t & > & \frac{20}{13-\sqrt{29}}\approx2.63\mbox{ hrs after 8AM}
\end{eqnarray*}

So about $8+2.63=10.63\mbox{ AM}$ or $10:\left(.63\right)\cdot60=10:38\mbox{ AM}$\vfill
\end{enumerate}
\newpage\textbf{If William gets a head start:}
\begin{enumerate}[resume]
\item William, Nancy, and Edgar met at the lodge at \uline{8 AM}, and William
began hiking at 2 mph. \uline{At 10 AM}, Nancy began hiking north
at 5 mph and Edgar went East on a three-wheeler at 12 mph. At what
time was the distance between Nancy and Edgar 20 miles greater than
than the distance between Nancy and William?

\vspace{.5cm}We need to solve the inequality 
\begin{equation}
\mbox{dist}\left(N,E\right)>\mbox{dist}\left(N,W\right)+20\label{eq:dist}
\end{equation}
Where $N=\mbox{Nancy's position}$, $E=\mbox{Edgar's position}$,
and $W=\mbox{William's position}$ as $\left(x,y\right)$ coordinates
with the cabin at the origin (point $(0,0)$). See the figure:

\begin{center}
\renewcommand{\xmax}{15}
\renewcommand{\xmin}{-8}
\renewcommand{\xstep}{0} %must be xmin+n
\renewcommand{\ymax}{8}
\renewcommand{\ymin}{0}
\renewcommand{\ystep}{0} %must be ymin+n
%%%%%%%
\begin{tikzpicture} 
\begin{axis}[height=3.5cm, 
	width=3.5cm, 
	axis lines=middle,
	minor x tick num={4},
	minor y tick num={4},
	grid=none,
	xtick={0,50,...,250},
	ytick={0,10,20,...,40},
	%axis line style={-latex}, 
	xmin=\xmin,xmax=\xmax, 
	ymin=\ymin,ymax=\ymax,
	%restrict y to domain=-1:10,
	%no markers, 
	xlabel={$$},ylabel={$$}, 
	%every axis x label/.append style={anchor=north}, 
	%every axis y label/.append style={anchor=east}
	]
	
	\addplot[color=red,mark=o,solid] coordinates {(0,5)} node[right] {$N$};
	\addplot[color=blue,mark=o,solid] coordinates {(12,0)} node[above] {$E$};
	\addplot[color=black,mark=o,solid] coordinates {(-6,0)} node[above] {$W$};

	\addplot[color=blue,dashed,mark=*,mark size=1pt] coordinates {(0,5) (12,0) };
	\addplot[color=red,dashed,mark=*,mark size=1pt] coordinates {(0,5) (-6,0) };

	
\end{axis}
\end{tikzpicture} 
\par\end{center}

$N=\left(0,5t\right)$, $E=\left(12t,0\right)$, and $W=\left(-2t-4,0\right)$.
William starts 2 hours earlier hence the $-2\cdot2=-4$ (the negative
denotes direction; it turns out to be unnecessary in this problem).
Each point is given with respect to time \uline{after 10 AM}. At $t=0$,
$N$ and $E$ they are at the cabin, $\left(0,0\right)$; $W$ is
4 miles west of the cabin, $\left(-4,0\right)$. They all move away
from the cabin according to their speed and direction. The picture
shows $t=1$.

$\mbox{dist}\left(\left(x_{1},y_{1}\right),\left(x_{2},y_{2}\right)\right)=\sqrt{\left(x_{1}-x_{2}\right)^{2}+\left(y_{1}-y_{2}\right)^{2}}$,
so:
\begin{itemize}
\item $\mbox{dist}\left(N,W\right)=\sqrt{\left(0-\left(-2t-4\right)\right)^{2}+\left(5t-0\right)^{2}}=\sqrt{4t^{2}+16t+16+25t^{2}}=\sqrt{29t^{2}+16t+4}$
\item $\mbox{dist}\left(N,E\right)=\sqrt{\left(0-12t\right)^{2}+\left(5t-0\right)^{2}}=\sqrt{144t^{2}+25t^{2}}=\sqrt{169t^{2}}=13t$
\end{itemize}
Plugging this into equation \ref{eq:dist}, we have:
\begin{eqnarray}
13t & > & \sqrt{29t^{2}+16t+4}+20\nonumber \\
13t-20 & > & \sqrt{29t^{2}+16t+4}\label{eq:ineq}\\
\left(13t-20\right)^{2} & > & \left(\sqrt{29t^{2}+16t+4}\right)^{2}\nonumber \\
169t^{2}-520t+400 & > & 29t^{2}+16t+4\nonumber \\
140t^{2}-536t+396 & > & 0\nonumber \\
4(t-1)(35t-99) & > & 0\nonumber 
\end{eqnarray}

So we need to solve the quadratic inequality, $4(t-1)(35t-99)=0$.
So $t=1$ or $t=\frac{99}{35}\approx2.829$. But $t=1\Rightarrow13(1)-20=-7$
which is not a solution for the inequality \ref{eq:ineq} above because
the radical cannot be less than 0. So then $t\approx2.829\mbox{ hrs}$. 

So the solution is approximately $10+2.829=12.83\mbox{ PM}$ or $12:50\mbox{ PM}.$
\end{enumerate}
\vspace{1cm}While a good problem, it is difficult because there are
so many little details to mess up!

Note that the negative used for $W$'s position was unnecessary because
it got squared. However if you wanted to find the distance between
$W$ and $E$ it would be needed. the distance between points could
have been found using Pythagoras\textquoteright s Theorem, but this
theorem is really just a distance formula. 
\end{document}
